\documentclass[11pt]{article}
\usepackage[T1]{fontenc}
\usepackage{lmodern}
\usepackage[margin=1.05in]{geometry}
\usepackage{amsmath,amssymb,amsthm}
\usepackage{booktabs}
\usepackage{microtype}
\microtypesetup{expansion=false}
\usepackage{xcolor}
\definecolor{linknavy}{RGB}{25,55,125}
\usepackage[colorlinks=true,linkcolor=linknavy,citecolor=linknavy,urlcolor=linknavy]{hyperref}
\hypersetup{
  pdfauthor={Felipe Santibanez-Leal},
  pdftitle={Curvilinear Geometry, Primary Structure, and Groebner Degenerations of Conductor Fiber Cones in a Huneke--Wiegand Family},
  pdfsubject={Truncated monomial parametrizations, primary decomposition, curvilinear geometry, and explicit Groebner bases},
  pdfkeywords={Huneke-Wiegand conjecture, numerical semigroups, conductor ideals, fiber cones, primary ideals, curvilinear schemes, fat points, Groebner bases, monomial degenerations, Kahler differentials}
}
\usepackage{orcidlink}

\newtheorem{theorem}{Theorem}[section]
\newtheorem{proposition}[theorem]{Proposition}
\newtheorem{corollary}[theorem]{Corollary}
\newtheorem{lemma}[theorem]{Lemma}
\theoremstyle{definition}
\newtheorem{definition}[theorem]{Definition}
\theoremstyle{remark}
\newtheorem{remark}[theorem]{Remark}

\newcommand{\N}{\mathbb{N}}
\newcommand{\kfield}{\Bbbk}
\newcommand{\Cfiber}{\mathcal{C}}
\newcommand{\Fcone}{\mathcal{F}}

\title{Curvilinear Geometry, Primary Structure, and Gr\"obner Degenerations\
of Conductor Fiber Cones in a Huneke--Wiegand Family}
\author{Felipe Santiba\~nez-Leal\,\orcidlink{0000-0002-0150-3246}\\
\small CAOS Research programme, \texttt{github.com/fsantibanezleal/CAOS\_RESEARCH}}
\date{2026-08-18\\[2pt]\normalsize Version 0.02}

\begin{document}
\maketitle
\begin{center}\small
\fbox{\parbox{0.92\textwidth}{\centering
\textbf{PREPRINT} \textperiodcentered\ not peer reviewed \textperiodcentered\ CAOS Research programme\\
Companion paper in the Huneke--Wiegand extension record \textperiodcentered\ version 0.02 \textperiodcentered\ 2026-08-18\\
Version DOI reserved only after complete pre-upload validation \textperiodcentered\ licensed CC BY 4.0\\
Concept DOI: \href{https://doi.org/10.5281/zenodo.21997377}{10.5281/zenodo.21997377}\\
Sources and machine records: \texttt{github.com/fsantibanezleal/CAOS\_RESEARCH}\\[2pt]
\textbf{Provenance.} Son Pham discovered the first public counterexample addressed by the wider
programme, and Professor Craig Huneke independently verified it. This companion paper concerns a
separate CAOS infinite family and the projective geometry of its conductor fiber cones.
}}
\end{center}

\begin{abstract}
For every integer $p\geq4$, a previously constructed symmetric numerical semigroup ring $R_p$
and rigid two-generated monomial ideal determine a conductor ideal $T_p$ whose special fiber
$\Cfiber_p=\Fcone(T_p)$ is one-dimensional Cohen--Macaulay, generated in degree one by $10p$
monomials, and has multiplicity $24p$. We prove that this fiber cone admits the explicit
standard-graded parametrization
\[
 \Cfiber_p\cong
 \kfield[xy^a:a\in G_p]\subseteq \kfield[x,y]/(y^{24p}),
 \qquad \deg(xy^a)=1,
\]
where $G_p$ is the exact degree-one offset set. If
$P_p=\kfield[X_a:a\in G_p]$ and $J_p$ is the defining ideal, then
\[
 \sqrt{J_p}=L_p=(X_a:a>0),\qquad J_p\text{ is }L_p\text{-primary}.
\]
Thus the complete primary decomposition has one component. The nilradical of $\Cfiber_p$ has
sharp nilpotency index $24p$: its $(24p-1)$st power is nonzero, while its $24p$th power vanishes.
Dehomogenization at $X_0$ gives
\[
 \Cfiber_p/(X_0-1)\cong\kfield[y]/(y^{24p}).
\]
Consequently $\operatorname{Proj}(\Cfiber_p)$ is a saturated length-$24p$ curvilinear fat point
with one-dimensional Zariski tangent space. It is locally Gorenstein, although its homogeneous
coordinate ring has Cohen--Macaulay type $10p+1$ and is neither level nor Gorenstein. We also
compute its affine K\"ahler differential module, including the exact characteristic split. The
proof is deductive from frozen offset-basis and Cohen--Macaulay input theorems. We further give the
complete reduced Gr\"obner basis for graded reverse lexicographic order with $X_0$ last. Its degree
profile is
\[
 (50p^2-17p,\;5p-1,\;p-2)
\]
in degrees two, three, and four, with no later elements. No minimal leading monomial contains
$X_0$, yielding a flat Cohen--Macaulay monomial degeneration. Exact campaigns, an independently
encoded clique audit, and an all-parameter Presburger boundary certificate validate the formulas
and adversarial controls but do not replace the deductive arguments.
\end{abstract}

\section{Introduction}

Fiber cones encode the number and relations of generators of powers of an ideal. Even for
monomial ideals, their reducedness and projective geometry can vary sharply. Explicit defining
ideals are known for several convex, concave, symmetric, and low-generator classes
\cite{HerzogQureshiSaem2017,HerzogZhu2019}; in particular, nonradical monomial fiber cones occur
naturally. Curvilinear zero-dimensional schemes provide a complementary geometric language: their
local maximal ideals are generated by one element, and their differential modules give effective
signatures of this property \cite{KreuzerLinhLong2023}.

This paper brings those two viewpoints together for the conductor fiber cones attached to an
explicit infinite family in the Huneke--Wiegand extension programme. The family was constructed
in \cite{SantibanezLeal2026} after Son Pham's discovery of a first public counterexample
\cite{Pham2026}; the discovery priority and the present contribution are distinct. The same
earlier work determines the conductor fiber cone, its complete degreewise offset bases, its
Cohen--Macaulay type, its defining ideal, and exact edges of its graded resolution. Here we use
that information from a different angle. The presence of offsets $0$ and $1$, combined with a
hard truncation at $24p$, makes the entire projective scheme one-variable after dehomogenization.

The main result is stronger than a radical calculation. It gives a concrete truncated monomial
model, a complete one-component primary decomposition, the exact nilpotency exponent, and the
scheme-theoretic affine chart. It also separates two notions that are easy to conflate: the
projective scheme is locally Gorenstein, while its given homogeneous embedding is not
arithmetically Gorenstein.

Throughout, derived theorem statements are marked [D] by this declaration, while finite exact
campaigns and solver outcomes are machine-verified [MV]. No [MV] record is used as an
extrapolation principle.

\section{The conductor family and frozen input}

Fix $p\geq4$ and put $s=6p$ and $q=4s=24p$. Define residue sets in $[0,s-1]$ by
\begin{align*}
A_p&=[0,p]\cup[3p,4p-2],\\
B_p&=\bigl([p+1,3p-1]\setminus\{2p-1\}\bigr)
      \cup\{4p\}\cup[5p-1,6p-1],\\
D_p&=[0,2p]\cup[3p,5p-2].
\end{align*}
The numerical semigroup $\Gamma_p$ contains
\[
\{0\},\quad 4s+A_p,\quad [5s,6s-1],\quad 6s+B_p,\quad 8s+D_p,
\quad [9s,13s-2],\quad [13s,\infty),
\]
and no other nonnegative integers. The earlier family theorem proves that $\Gamma_p$ is symmetric,
with multiplicity $24p$ and Frobenius number $78p-1$, and that
$I_p=(t^{24p},t^{30p})$ is a nonprincipal rigid ideal in the localization of
$\kfield[t^{\Gamma_p}]$ at its positive-degree maximal ideal \cite{SantibanezLeal2026}.

Let $E_p=\operatorname{End}_{R_p}(I_p)$ and let $T_p=R_p:E_p$ be the conductor. The object studied
here is its special fiber
\[
\Cfiber_p=\Fcone(T_p)=\bigoplus_{n\geq0}T_p^n/\mathfrak m_pT_p^n.
\]
We recall the exact input in the form needed below.

\begin{proposition}[Frozen offset-basis input]\label{prop:input}
There is a presentation
\[
 P_p=\kfield[X_a:a\in G_p]\longrightarrow\Cfiber_p,
 \qquad X_a\longmapsto[t^{q+a}],
\]
where $|G_p|=10p$ and
\begin{align*}
G_p={}&\{0\}\cup[1,p]\cup[3p,4p-2]\cup[6p,8p-2]\cup[8p,10p-2]
\cup\{10p\}\\
&\cup[11p-1,12p-1]\cup[13p+1,14p-2]\cup[14p,15p-1]
\cup\{16p\}\cup[17p-1,18p-1].
\end{align*}
If $J_p$ is the kernel, a degree-$n$ monomial is zero precisely when its total offset is absent
from $E_n$, and two nonzero monomials have the same class precisely when their total offsets
agree. The cumulative bases are
\begin{align*}
E_1&=G_p,\\
E_2&=[0,2p]\cup[3p,5p-2]\cup[6p,q-1],\\
E_3&=[0,q-1]\setminus\{6p-1\},\\
E_n&=[0,q-1]\qquad(n\geq4).
\end{align*}
Moreover, $\Cfiber_p$ is a one-dimensional Cohen--Macaulay standard graded algebra, free of rank
$q$ over $\kfield[X_0]$, and has Cohen--Macaulay type $10p+1$. Its Artinian reduction by $X_0$
has socle in two degrees, so $\Cfiber_p$ is neither level nor Gorenstein.
\end{proposition}

\begin{remark}
Proposition~\ref{prop:input} packages the previously proved conductor, fiber-cone, presentation,
and offset-congruence theorems. The present paper does not infer it from the finite campaign below.
Its exact all-parameter component step retains the solver/encoding boundary disclosed in
\cite{SantibanezLeal2026}.
\end{remark}

\section{The truncated monomial model}

Assign standard degree one to each monomial $xy^a$ and define
\[
 \mathcal A_p=\kfield[xy^a:a\in G_p]\subseteq\kfield[x,y]/(y^q).
\]

\begin{theorem}[Truncated parametrization]\label{thm:parametrization}
For every $p\geq4$, the homomorphism
\[
 P_p\longrightarrow\mathcal A_p,
 \qquad X_a\longmapsto xy^a,
\]
induces a standard-graded isomorphism $\Cfiber_p\cong\mathcal A_p$.
\end{theorem}

\begin{proof}
A degree-$n$ monomial with total offset $b$ maps to $x^ny^b$. This is zero exactly when
$b\geq q$, and distinct surviving offsets give linearly independent monomials. By
Proposition~\ref{prop:input}, the corresponding class in $(\Cfiber_p)_n$ is zero precisely when
the offset is absent from $E_n$, and otherwise depends only on the offset. The achieved surviving
offsets are exactly $E_n$ in both quotients. Hence the two kernels agree in every graded degree.
\end{proof}

\begin{corollary}[Exact dehomogenization]\label{cor:dehom}
There is an isomorphism
\[
 \Cfiber_p/(X_0-1)\cong\kfield[y]/(y^q),
 \qquad X_a\longmapsto y^a.
\]
\end{corollary}

\begin{proof}
The offsets $0$ and $1$ belong to $G_p$. For every positive $a\in G_p$, the two degree-$a$
monomials $X_0^{a-1}X_a$ and $X_1^a$ have the same surviving offset $a<q$, so
\[
 X_0^{a-1}X_a=X_1^a.
\]
After setting $X_0=1$, every generator is a power of $X_1$, and $X_1^q=0$. This gives a
surjection from $\kfield[y]/(y^q)$. Freeness of rank $q$ over $\kfield[X_0]$ shows that the target
also has vector-space dimension $q$, hence the surjection is an isomorphism.
\end{proof}

\section{Primary decomposition and sharp nilpotence}

Let
\[
 L_p=(X_a:a>0)\subseteq P_p,
 \qquad N_p=L_p/J_p\subseteq\Cfiber_p.
\]

\begin{theorem}[One primary component]\label{thm:primary}
For every $p\geq4$,
\[
 \sqrt{J_p}=L_p,\qquad J_p\text{ is }L_p\text{-primary}.
\]
Thus $J_p$ itself is the complete primary decomposition. The nilradical $N_p$ has exact
nilpotency index $q=24p$:
\[
 N_p^{q-1}\neq0,\qquad N_p^q=0.
\]
\end{theorem}

\begin{proof}
The quotient $\Cfiber_p/N_p$ is $\kfield[X_0]$, so it is reduced. Every product of $q$ generators
with positive offset has total offset at least $q$ and vanishes by
Theorem~\ref{thm:parametrization}. Hence $N_p$ is nilpotent and equals the nilradical. It is the
unique minimal prime. The element $X_1^{q-1}$ maps to $x^{q-1}y^{q-1}$ and is nonzero, proving
the sharp exponent.

The unique-minimal-prime statement alone would not prove primaryness. Here the
Cohen--Macaulay hypothesis is essential: a one-dimensional Cohen--Macaulay ring has no embedded
associated primes. Therefore $\operatorname{Ass}(\Cfiber_p)=\{N_p\}$, so the zero ideal is
$N_p$-primary. Taking its inverse image in $P_p$ proves the assertion about $J_p$.
\end{proof}

\begin{remark}[Why Cohen--Macaulayness is load-bearing]
The ring $\kfield[u,v]/(u^2,uv)$ has a unique minimal prime but also an embedded associated prime;
its zero ideal is not primary. This is an explicit control against replacing the associated-prime
argument by an invalid radical-only inference.
\end{remark}

\section{Curvilinear geometry and Gorenstein separation}

\begin{theorem}[The projective fat point]\label{thm:geometry}
The ideal $J_p$ is saturated for the irrelevant ideal. The projective scheme
$\operatorname{Proj}(\Cfiber_p)$ is supported at $[1:0:\cdots:0]$ and is a length-$q$
curvilinear fat point. Its Zariski tangent space has dimension one. The scheme is locally
Gorenstein, while the homogeneous coordinate ring $\Cfiber_p$ is not level and not Gorenstein.
\end{theorem}

\begin{proof}
Because $\Cfiber_p$ is one-dimensional Cohen--Macaulay, its zeroth local cohomology for the
irrelevant ideal vanishes. The identity
$J_p^{\mathrm{sat}}/J_p=H^0_{P_{p,+}}(\Cfiber_p)$ gives saturation. Theorem~\ref{thm:primary}
shows that the support is the single coordinate point and hence lies entirely in $D_+(X_0)$.
Corollary~\ref{cor:dehom} identifies its affine coordinate ring with
$\kfield[y]/(y^q)$. This has length $q$ and principal maximal ideal $(y)$, so the scheme is
curvilinear and its tangent space $(y)/(y^2)$ is one-dimensional. Its socle is generated by
$y^{q-1}$, proving local Gorensteinness. Proposition~\ref{prop:input} gives type $10p+1$ and
nonlevelness for the homogeneous coordinate ring, proving the arithmetic contrast.
\end{proof}

\begin{remark}
This example makes the embedding distinction concrete:
\[
 \text{locally Gorenstein}\quad\not\Rightarrow\quad\text{arithmetically Gorenstein}.
\]
The local scheme is always the one-variable truncation, while the standard grading retains the
nonlevel offset-layer structure.
\end{remark}

\section{K\"ahler differentials}

The curvilinear chart has a short differential fingerprint which also records the ground-field
characteristic.

\begin{proposition}[Differential signature]\label{prop:differentials}
Let $B_p=\Cfiber_p/(X_0-1)$. Then
\[
 \Omega^1_{B_p/\kfield}
 \cong \bigl(B_p/(q y^{q-1})\bigr)\,dy,
 \qquad \bigwedge
 ^m\Omega^1_{B_p/\kfield}=0\quad(m\geq2).
\]
Consequently,
\[
 \dim_{\kfield}\Omega^1_{B_p/\kfield}=
 \begin{cases}
 q-1,&\operatorname{char}(\kfield)\nmid q,\\
 q,&\operatorname{char}(\kfield)\mid q.
 \end{cases}
\]
The cotangent space at the closed point has dimension one in every characteristic.
\end{proposition}

\begin{proof}
Apply the conormal sequence to $B_p=\kfield[y]/(y^q)$. Differentiating the defining equation gives
$q y^{q-1}dy=0$, which proves the presentation and the characteristic split. The differential
module is cyclic, so every exterior power of degree at least two vanishes. Tensoring with
$B_p/(y)$ leaves the single generator $dy$.
\end{proof}

This agrees with the general differential characterization of weakly curvilinear schemes in
\cite{KreuzerLinhLong2023}, while the exact truncation and characteristic formula here come from
the conductor-family parametrization.

\section{The canonical grevlex degeneration}

Order the variables by decreasing offset and use graded reverse lexicographic order $\prec$; thus
$X_a\succ X_b$ exactly when $a>b$, and $X_0$ is last. If a monomial is written as a
nondecreasing tuple of offsets, then the $\prec$-smallest monomial in a fixed degree and offset
fiber is the lexicographically smallest tuple. Write it as $N_{n,b}$.
This last-parameter order and the use of finite standard-monomial counts are consistent with
Gr\"obner methods for projective monomial curves and associated graded semigroup algebras
\cite{BhardwajChauJavadekar2025,SahaSenguptaSrivastava2022}; the staircase and formulas below are
specific to the present truncated conductor family.

\begin{lemma}[Canonical staircase]\label{lem:staircase}
In each surviving bidegree $(n,b)$, the unique standard monomial of
$P_p/\operatorname{in}_{\prec}(J_p)$ is $N_{n,b}$. A monomial is a minimal generator of the
initial ideal exactly when it is noncanonical and every one-variable divisor is canonical.
\end{lemma}

\begin{proof}
Theorem~\ref{thm:parametrization} shows that two nonzero degree-$n$ monomials agree precisely when
their total offsets agree, and that all monomials of offset at least $q$ vanish. Every surviving
fiber therefore contributes its unique $\prec$-smallest monomial to the standard basis. The last
assertion is the minimal-boundary criterion for a monomial ideal.
\end{proof}

\begin{theorem}[Explicit reduced Gr\"obner basis]\label{thm:groebner}
For every $p\geq4$, the reduced Gr\"obner basis of $J_p$ has degree profile
\[
 \bigl(50p^2-17p,\;5p-1,\;p-2\bigr)
\]
in degrees two, three, and four. There are no basis elements in degree at least five.

The quadratic part contains one relation for every noncanonical quadratic monomial $M$ of total
offset $b$:
\[
 M-N_{2,b}\quad(b<q),\qquad M\quad(b\geq q).
\]
Among these, $(77p^2-49p+2)/2$ are binomials and $(23p^2+15p-2)/2$ are monomial
zero-relations.

The cubic leading monomials are
\begin{align*}
A_i&=X_iX_pX_{12p-1} &&(1\leq i\leq p),\\
B_i&=X_iX_pX_{15p-1} &&(1\leq i\leq p),\\
C_i&=X_iX_pX_{18p-1} &&(1\leq i\leq p),\\
D_i&=X_iX_{4p-2}X_{16p} &&(1\leq i\leq p-2),\\
E_i&=X_iX_{4p-2}X_{18p-1} &&(1\leq i\leq p),
\end{align*}
together with $X_p^3$. Their reduced tails are
\begin{align*}
A_1&\longmapsto X_0X_{3p}X_{10p},&
A_i&\longmapsto X_0^2X_{13p+i-1} &&(2\leq i\leq p-1),\\
A_p&\longmapsto X_0X_1X_{14p-2},&
B_i&\longmapsto X_0X_{i-1}X_{16p} &&(1\leq i\leq p-1),\\
B_p&\longmapsto X_0^2X_{17p-1},&
C_i&\longmapsto X_0X_{3p+i-1}X_{16p} &&(1\leq i\leq p-1),\\
C_p&\longmapsto X_0X_{3p}X_{17p-1},&
D_i&\longmapsto X_0X_{3p}X_{17p+i-2},\\
E_i&\longmapsto X_0X_{7p+i-2}X_{15p-1} &&(i=1,2),\\
E_i&\longmapsto X_0X_{6p+i-3}X_{16p} &&(3\leq i\leq p),\\
X_p^3&\longmapsto X_0^2X_{3p}.&&
\end{align*}
Finally, the quartics are
\[
 X_iX_p^2X_{4p-2}-X_0^3X_{6p+i-2},\qquad 2\leq i\leq p-1.
\]
No minimal leading monomial is divisible by $X_0$.
\end{theorem}

\begin{proof}
There are $\binom{10p+1}{2}=50p^2+5p$ quadratic monomials and $|E_2|=22p$ canonical
ones, giving $50p^2-17p$ quadratic boundary elements. Summing the unordered pairs from the eleven
blocks of $G_p$ whose offset is at least $q$ gives $(23p^2+15p-2)/2$ zero-relations; subtraction
gives the stated binomial count.

For the nonquadratic boundary, Lemma~\ref{lem:staircase} turns completeness into Presburger
arithmetic over the eleven interval blocks of $G_p$. The exact certificate proves, for all
integers $p\geq4$, that the five indexed cubic groups plus $X_p^3$ are both sound and complete,
that the displayed quartic family is sound and complete, and that all twelve endpoint/interior
tail cases are canonical and have the correct offset. Its 16 negated obligations are UNSAT. The
family sizes give $p+p+p+(p-2)+p+1=5p-1$ cubics and $p-2$ quartics.

It remains to close the infinite degree tail without solver extrapolation. Since $E_4=[0,q-1]$
and $E_n=[0,q-1]$ for $n\geq4$, the last-variable order forces
\[
 N_{n,b}=X_0^{n-4}N_{4,b}\qquad(n\geq4,\ 0\leq b<q).
\]
The staircase therefore stabilizes in degree four, so no later minimal boundary can occur. The
listed polynomials are monic with distinct minimal leading monomials and canonical tails; hence
they form the reduced Gr\"obner basis.
\end{proof}

\begin{corollary}[Flat Cohen--Macaulay degeneration]\label{cor:degeneration}
The initial algebra is Cohen--Macaulay, $X_0$ is regular on it, and its Artinian reduction has
Hilbert function
\[
 (1,\;10p-1,\;12p,\;2p-1,\;1).
\]
The defining ideal has relation type three, while this reduced Gr\"obner basis reaches degree
four.
\end{corollary}

\begin{proof}
No minimal initial generator contains $X_0$, so multiplication by $X_0$ is injective on standard
monomials. The successive offset layers from Proposition~\ref{prop:input} have the displayed
sizes. Their sum is $24p$, and adjoining powers of $X_0$ recovers the stabilized staircase.
Relation type three was established by the earlier minimal-presentation theorem, whereas
Theorem~\ref{thm:groebner} has $p-2>0$ quartics.
\end{proof}

\section{Exact verification record}

The symbolic proofs above use the frozen results in Proposition~\ref{prop:input}. Computation is
used to validate their translation into the new parametrization and to exercise negative controls.
The primary-structure campaign computes truncated sumsets directly from $G_p$ for every
$p=4,\ldots,300$. Its audit starts from the five disjoint Artinian layers. The Gr\"obner campaign
independently checks all declared families, canonical tails, divisor minimality, Hilbert values,
and controls at the same 297 parameters. A separately encoded clique audit enumerates canonical
pairs and reconstructs factorization complexes through degree six at
$p=4,5,6,17,73,151,300$.

\begin{center}
\begin{tabular}{@{}lll@{}}
\toprule
record & coverage & SHA-256 aggregate\\
\midrule
primary campaign & all $p=4,\ldots,300$ & \texttt{f3373f4f58287fd3...f9d6b39b6}\\
primary audit & all rows; six detailed rebuilds & \texttt{84c00be8ff64002...f5b6c8f143e}\\
Gr\"obner campaign & all $p=4,\ldots,300$ & \texttt{63af8f734afc8c05...9e4365a218}\\
clique audit & seven detailed rebuilds & \texttt{401c4807cc0a29a6...586766e41373}\\
Presburger boundary & all integers $p\geq4$ & \texttt{10c66bbcaa56108f...e046f9dd5fa}\\
\bottomrule
\end{tabular}
\end{center}

The controls delete offset $1$, perturb the truncation and nilpotency exponents, corrupt the
degree-three hole, omit a radical coordinate, conflate local and arithmetic Gorensteinness, erase
the differential characteristic split, and remove the no-embedded-primes premise. Every mutation
is rejected. Exact scripts, artifacts, premise hashes, proofs, and verdicts are archived as
\texttt{EXP-025-curvilinear-primary-structure} and \texttt{EXP-026-grevlex-staircase} in the
public repository.

The finite sweep does not prove the theorem for all $p$. The infinite conclusion is supplied by
Theorems~\ref{thm:parametrization}--\ref{thm:geometry} and
Proposition~\ref{prop:differentials}. A defect in a frozen input theorem would propagate even if
both implementations agreed. The Presburger boundary certificate uses Z3 4.16.0 and does not
emit independently checked UNSAT proof objects; the canonical-staircase and stabilization
arguments remain separate deductive steps.

\section{Scope and further questions}

The result concerns the conductor fiber cones of the explicit family recalled in Section~2. It
does not assert that arbitrary monomial fiber cones are primary, curvilinear, Cohen--Macaulay, or
locally Gorenstein. It does not strengthen the underlying Huneke--Wiegand counterexample theorem
for arbitrary modules or arbitrary one-dimensional Gorenstein domains.

The primary decomposition and the reduced Gr\"obner basis are now complete. One major algebraic
question remains open for this family:
\begin{enumerate}
  \item Can the unresolved interior of the graded Betti table be expressed by closed formulas,
  or does it exhibit parameter-dependent combinatorial phases?
\end{enumerate}
The monomial degeneration may help by replacing value-semigroup products with the combinatorics
of one explicit finite staircase.

\section*{Acknowledgments}

The author thanks Son Pham for making the first counterexample and its verification materials
public, and Craig Huneke for the independent verification recorded with that release. Neither is
an author of this paper, and this acknowledgment does not imply endorsement of the present
theorem. All authorship and accountability for this manuscript remain with the named human author.

{\footnotesize
\begin{thebibliography}{99}
\setlength{\itemsep}{0.25em}
\setlength{\parsep}{0pt}

\bibitem{HunekeWiegand1994}
C. Huneke and R. Wiegand,
\emph{Tensor products of modules and the rigidity of Tor},
Math. Ann. 299 (1994), 449--476.
\href{https://doi.org/10.1007/BF01459793}{doi:10.1007/BF01459793}.

\bibitem{Pham2026}
S. Pham,
\emph{Huneke--Wiegand candidate verification}, public software and certificate repository (2026).
\url{https://github.com/sonpham-org/huneke-wiegand-candidate-verification}.

\bibitem{SantibanezLeal2026}
F. Santiba\~nez-Leal,
\emph{Frobenius minimality, minimum-layer uniqueness, and an infinite family of
Huneke--Wiegand counterexamples}, CAOS Research preprint, version 0.13 (2026).
\href{https://doi.org/10.5281/zenodo.21995498}{doi:10.5281/zenodo.21995498}.

\bibitem{HerzogQureshiSaem2017}
J. Herzog, A. A. Qureshi, and M. Mohammadi Saem,
\emph{The fiber cone of a monomial ideal in two variables}, arXiv:1711.08775 (2017).
\url{https://arxiv.org/abs/1711.08775}.

\bibitem{HerzogZhu2019}
J. Herzog and G. Zhu,
\emph{On the fiber cone of monomial ideals}, arXiv:1904.04988 (2019).
\url{https://arxiv.org/abs/1904.04988}.

\bibitem{KreuzerLinhLong2023}
M. Kreuzer, T. N. K. Linh, and L. N. Long,
\emph{Differential theory of zero-dimensional schemes}, arXiv:2302.11903 (2023).
\url{https://arxiv.org/abs/2302.11903}.

\bibitem{BhardwajChauJavadekar2025}
O. P. Bhardwaj, T. Chau, and O. Javadekar,
\emph{Projective monomial curves associated to numerical semigroups with multiplicity $e$,
width $e-1$, and embedding dimension $e-2$}, arXiv:2511.06482 (2025).
\url{https://arxiv.org/abs/2511.06482}.

\bibitem{SahaSenguptaSrivastava2022}
J. Saha, I. Sengupta, and P. Srivastava,
\emph{On the associated graded ring of semigroup algebras}, arXiv:2210.07520 (2022).
\url{https://arxiv.org/abs/2210.07520}.

\end{thebibliography}
}

\end{document}
